\section{Monetary policy as deployed}
\label{sec:monetary-policy}

\Flux issues its native asset on a schedule fixed once, in \fluxd's consensus code, and not
adjusted by discretion afterward. Two different mechanisms have produced that schedule over the
chain's life: a Bitcoin-style halving curve during the proof-of-work era, and, since block height
$2{,}020{,}000$, a fixed initial block reward that decays by a constant fraction once a year under
Proof of Node. Neither curve reaches zero on any realistic horizon, and at least one of them was
adjusted by hand along the way rather than left to run the formula a reader might expect. This
section states the exact, currently-enforced arithmetic for both eras in the notation of
notation.tex, cross-checks the resulting figures against the network's own public reconciliation
of issued supply, and states plainly what the schedule implies for the long run: issuance never
stops, and only its rate settles to a small constant. Every subsequent economic claim in this
paper is denominated against the numbers derived here.

All amounts below are computed as integers over satoshi, $\mathrm{COIN} = 100{,}000{,}000\sat$ per
$\flux$ \shipped \srccode{fluxd/\allowbreak{}src/\allowbreak{}amount.h}{17}, matching \fluxd's \texttt{CAmount}
(\texttt{int64}) representation. ``$\lfloor\cdot\rfloor$'' denotes truncating (floor) integer
division throughout, exactly as C++ integer division behaves for non-negative operands
\srcdoc{flux-roadmap/supply/emission\_model.py:196-200, citing fluxd}. This distinction is
not cosmetic: the PoN reduction below is an $r$-fold sequence of \emph{independent} truncations,
not a single closed-form power, and the two are not interchangeable.

\subsection{The subsidy function}

\begin{definition}[Legacy proof-of-work subsidy]
\label{def:pow-subsidy}
For $0 \le \hgt < \hpon$,
$$
\subsidypow(\hgt) \;\defeq\;
\begin{cases}
0 & \hgt = 0 \\[4pt]
\premine = 13{,}020{,}000\flux & \hgt = 1 \\[4pt]
\floorfrac{150\,\mathrm{COIN}}{5000}\cdot \hgt & 2 \le \hgt < 2500 \\[6pt]
\floorfrac{150\,\mathrm{COIN}}{5000}\cdot (\hgt+1) & 2500 \le \hgt < 5000 \\[6pt]
\left\lfloor \dfrac{150\,\mathrm{COIN}}{2^{\,\min\left(\floorfrac{\hgt-2500}{655{,}350},\,2\right)}} \right\rfloor & 5000 \le \hgt < \hpon
\end{cases}
$$
\end{definition}
\shipped \srccode{fluxd/\allowbreak{}src/\allowbreak{}main.cpp}{2818-2852}, constants \srccode{fluxd/\allowbreak{}src/\allowbreak{}chainparams.cpp}{90-91}.
This is the literal transcription of \texttt{Get\allowbreak{}Block\allowbreak{}Subsidy}'s legacy branch as currently compiled
\shipped \srccode{flux-roadmap/\allowbreak{}supply/\allowbreak{}emission\_model.py}{207-218}. The chain's actual historical
payout for heights $1$--$4{,}999$ differs from this formula by one block index (an older binary
paid the premine at height $2$, not $1$, and shifted the slow-start ramp accordingly)
\shipped \srccode{flux-roadmap/\allowbreak{}supply/\allowbreak{}emission\_model.py}{220-229}; that discrepancy is the source
of the persistent offset discussed in \S\ref{sec:mp-discrepancy} below. The code's separate
$\text{halvings}\ge 64 \Rightarrow 0$ branch is never reached before $\hpon$ and so never fires on
the deployed chain \shipped \srccode{fluxd/\allowbreak{}src/\allowbreak{}main.cpp}{2839-2840}.

\begin{definition}[Proof-of-Node subsidy]
\label{def:pon-subsidy}
For $\hgt \ge \hpon$, let
\begin{gather*}
r(\hgt) \;\defeq\; \min\!\left(\floorfrac{\hgt-\hpon}{\redint},\ \redmax\right),
\qquad
s_0 \defeq 14\cdot\mathrm{COIN} = 1{,}400{,}000{,}000\sat, \\
s_{k+1} \defeq \left\lfloor s_k \cdot \redfac \right\rfloor,\ \ \redfac = \frac{9}{10},
\end{gather*}
for $k = 0,\dots,\redmax-1$. Then
$$
\subsidypon(\hgt) \;\defeq\; s_{r(\hgt)}.
$$
\end{definition}
\shipped \srccode{fluxd/\allowbreak{}src/\allowbreak{}main.cpp}{2799-2816}; $\hpon = 2{,}020{,}000$
\shipped \srccode{fluxd/\allowbreak{}src/\allowbreak{}chainparams.cpp}{153}; $\redint = 1{,}051{,}200$ blocks
($365\times24\times60\times2$, one year at the 30-second PoN spacing)
\shipped \srccode{fluxd/\allowbreak{}src/\allowbreak{}chainparams.cpp}{157}; $\redmax = 20$
\shipped \srccode{fluxd/\allowbreak{}src/\allowbreak{}chainparams.cpp}{158}. The combined block subsidy is
$\subsidy(\hgt) = \subsidypow(\hgt)$ for $\hgt < \hpon$ and $\subsidypon(\hgt)$ otherwise
\shipped \srccode{fluxd/\allowbreak{}src/\allowbreak{}main.cpp}{2796-2798}.

\textbf{Integer semantics.} $s_{r}$ is $r$ \emph{independently truncated} multiplications, not a
closed form: $s_r \neq \left\lfloor s_0 \cdot (9/10)^{r} \right\rfloor$ in general, because each
step's rounding error compounds on top of the previous step's rather than being computed once
against the original $s_0$ \shipped \srccode{fluxd/\allowbreak{}src/\allowbreak{}main.cpp}{2811-2813}. At $r=20$ the two
expressions already differ at the eight printed decimal places used below: the iterated value is
$1.70207313\flux$ against $1.70207316\flux$ for the closed form, and the two first diverge at $r=10$
\shipped \srccode{flux-roadmap/\allowbreak{}supply/\allowbreak{}emission\_model.py}{195-204}; the paper
states the iterated definition as the one \fluxd actually computes.

\subsection{The frozen halving}
\label{sec:mp-frozen-halving}

The $\min(\cdot,2)$ term in Definition~\ref{def:pow-subsidy} is not an idealization; it is a
hand-frozen consensus rule. \texttt{Get\allowbreak{}Block\allowbreak{}Subsidy}'s legacy branch reads
\texttt{if (halvings >= 2) \{ nSubsidy >>= 2; return nSubsidy; \}}, with an in-source comment
describing it as a ``temp solution while we work on the next blockchain improvement protocols''
\shipped \srccode{fluxd/\allowbreak{}src/\allowbreak{}main.cpp}{2842-2848}. Emission was therefore held at the second-halving
level, $37.5\flux$/block, rather than continuing to halve every $655{,}350$ blocks
\shipped \srccode{fluxd/\allowbreak{}src/\allowbreak{}chainparams.cpp}{91}. This was not merely a theoretical stopping point:
had halving continued, a third halving (to $18.75\flux$/block) would have taken effect at height
$1{,}968{,}550 = 2500 + 3\times655{,}350$, which is before $\hpon = 2{,}020{,}000$; instead, the
frozen floor held the reward at $37.5\flux$/block for the roughly $51{,}450$ blocks between that
point and PoN activation, at which point the entire legacy path was replaced by
Definition~\ref{def:pon-subsidy} \shipped \srccode{fluxd/\allowbreak{}src/\allowbreak{}main.cpp}{2845-2848}.

\subsection{Tier partition and the dev-fund remainder}
\label{sec:mp-tier-devfund}

\begin{definition}[PoN tier reward]
For tier $X \in \tiers = \{\tC,\tN,\tS\}$ and $\hgt \ge \hpon$,
$$
\mathrm{reward}_X(\hgt) \;\defeq\; \left\lfloor \subsidypon(\hgt) \cdot \frac{\tbase{X}}{14\flux} \right\rfloor,
$$
with $\tbase{\tC} = 1.0\flux$, $\tbase{\tN} = 3.5\flux$, $\tbase{\tS} = 9.0\flux$ --- ratios chosen
so that two Nimbus earn less than one Stratus while three earn more, and three Cumulus less than one
Nimbus while four earn more \citep{fluxrewardmodel} ---
\shipped \srccode{fluxd/\allowbreak{}src/\allowbreak{}main.cpp}{2892-2895}, applied against the current, possibly-reduced
$\subsidypon(\hgt)$, not against a value fixed at PoN activation
\shipped \srccode{fluxd/\allowbreak{}src/\allowbreak{}main.cpp}{2886-2939}.
\end{definition}

The $14\flux$ denominator is coded as \texttt{PON\_INITIAL\_TOTAL}, a separate local constant in
\texttt{Get\allowbreak{}Fluxnode\allowbreak{}Subsidy}, rather than read from the consensus parameter
\texttt{n\allowbreak{}PONInitial\allowbreak{}Subsidy} that fixes $s_0$ in Definition~\ref{def:pon-subsidy}
\shipped \srccode{fluxd/\allowbreak{}src/\allowbreak{}main.cpp}{2892}. A future upgrade that changed
\texttt{n\allowbreak{}PONInitial\allowbreak{}Subsidy} without also editing this literal would silently desynchronize the
tier ratios from the actual base subsidy; this is a maintenance hazard in the deployed code, stated
here because it bears on what a future constant change could do without touching the tier logic at
all \shipped \srccode{fluxd/\allowbreak{}src/\allowbreak{}main.cpp}{2892}.

\begin{definition}[Dev-fund remainder]
$$
\devfund(\hgt) \;\defeq\; \subsidypon(\hgt) - \mathrm{reward}_{\tC}(\hgt) - \mathrm{reward}_{\tN}(\hgt) - \mathrm{reward}_{\tS}(\hgt), \qquad \hgt \ge \hpon.
$$
\end{definition}
\shipped \srccode{fluxd/\allowbreak{}src/\allowbreak{}main.cpp}{2877-2884} (\texttt{Get\allowbreak{}Min\allowbreak{}Dev\allowbreak{}Fund\allowbreak{}Amount}). At $r=0$, the
unreduced $14\flux$ subsidy, $\devfund = 0.5\flux = 50{,}000{,}000\sat$, i.e.\ $1/28 \approx 3.57\%$
of PoN emission. Because $\devfund$ is defined as a remainder of the \emph{current} subsidy rather
than a fixed floor, this share decays proportionally as $\subsidypon$ reduces
\shipped \srccode{flux-roadmap/\allowbreak{}supply/\allowbreak{}emission\_model.py}{358-368}.

Payment is consensus-enforced twice, independently. Loosely: any post-PoN coinbase lacking an
output of at least $\devfund(\hgt)$ to the hardcoded address is rejected with
\texttt{"tx-\allowbreak{}coinbase-\allowbreak{}missing-\allowbreak{}dev-\allowbreak{}funding"} \shipped \srccode{fluxd/\allowbreak{}src/\allowbreak{}main.cpp}{1231-1249}.
Strictly: \texttt{Fluxnode\allowbreak{}Cache::\allowbreak{}Check\allowbreak{}Fluxnode\allowbreak{}Payout} independently requires the coinbase value not
already claimed by a tier payout to be at least $\devfund(\hgt)$
\shipped \srccode{fluxd/\allowbreak{}src/\allowbreak{}fluxnode/\allowbreak{}fluxnode.cpp}{837-875}. The mainnet address is
\texttt{t3h\allowbreak{}Pu1YDe\allowbreak{}GUCp8m7BQCnn\allowbreak{}NUm\allowbreak{}RMJBa5Rady\allowbreak{}A} \shipped \srccode{fluxd/\allowbreak{}src/\allowbreak{}chainparams.cpp}{306}.

No source comment names this remainder a ``dev fund'' at the point where the tier ratios
themselves are defined; the routing is discoverable only by following
\texttt{Get\allowbreak{}Min\allowbreak{}Dev\allowbreak{}Fund\allowbreak{}Amount}'s call graph out to the two enforcement sites above, not from any
single comment at the site of the constants \shipped \srccode{fluxd/\allowbreak{}src/\allowbreak{}main.cpp}{2877-2884}. That
is a documentation gap worth stating plainly rather than leaving for a reader to reconstruct.

\subsection{One-off and recurring fund payouts}
\label{sec:mp-fund-payouts}

Three further payments are enforced directly in coinbase validation, by exact scriptPubKey and
amount match, independent of the tier/dev-fund split above
\shipped \srccode{fluxd/\allowbreak{}src/\allowbreak{}main.cpp}{1252-1304}: two one-off payments and one recurring schedule.

\begin{table}[H]
\centering
\caption{One-off and recurring consensus-enforced fund payouts, mainnet. All entries \shipped,
cited directly to \fluxd.}
\label{tab:fund-payouts}
\begin{tabular}{@{}lllll@{}}
\toprule
Payment & Address & Height / schedule & Amount & Provenance \\
\midrule
Exchange fund & \texttt{t3PMbb\allowbreak{}A5YBMrj\allowbreak{}SD3d\allowbreak{}D16SSd\allowbreak{}XKu\allowbreak{}Kovwmj6t\allowbreak{}S} & $\hgt = 836{,}274$, one-off &
$7{,}500{,}000\flux$ & \srccode{fluxd/\allowbreak{}src/\allowbreak{}chainparams.cpp}{293-294} \\
Foundation fund & \texttt{t3Xj\allowbreak{}YMBvwxn\allowbreak{}XVv9jqg4Cgok\allowbreak{}Z3f7k\allowbreak{}Ao\allowbreak{}XPQL8} & $\hgt = 836{,}994$, one-off &
$2{,}500{,}000\flux$ & \srccode{fluxd/\allowbreak{}src/\allowbreak{}chainparams.cpp}{297-298} \\
Swap pool & \texttt{t3Thb\allowbreak{}Wog\allowbreak{}Do\allowbreak{}Aj\allowbreak{}Gu\allowbreak{}S6DEnm\allowbreak{}N1GWJBRb\allowbreak{}Vj\allowbreak{}SUK4T} & from $\hgt = 837{,}714$, every
$21{,}600$ blocks, $\times 10$ & $22{,}000{,}000\flux$ per tranche ($220{,}000{,}000\flux$ total) &
\srccode{fluxd/\allowbreak{}src/\allowbreak{}chainparams.cpp}{301-304} \\
\bottomrule
\end{tabular}
\end{table}

\subsection{The supply function}
\label{sec:mp-supply}

\begin{definition}[Cumulative supply]
\label{def:supply}
\begin{multline*}
\supply(\hgt) \;=\; \mathrm{premine}(\hgt) + \mathrm{pow\_slow\_start}(\hgt)
  + \mathrm{pow}_{150}(\hgt) + \mathrm{pow}_{75}(\hgt) + \mathrm{pow}_{37.5}(\hgt) \\
  {}+ \mathrm{exch}(\hgt) + \mathrm{found}(\hgt) + \mathrm{swap}(\hgt) + \mathrm{pon}(\hgt),
\end{multline*}
with
\begin{align*}
\mathrm{premine}(\hgt) &= \premine \cdot \indic{\hgt \ge 2} \\
\mathrm{pow\_slow\_start}(\hgt) &= \textstyle\sum_{k=2}^{\min(\hgt,\,4999)} \subsidypow(k) \\
\mathrm{pow}_{150}(\hgt) &= 150\flux \cdot \big|\{k : 5000 \le k \le \min(\hgt,657{,}849)\}\big| \\
\mathrm{pow}_{75}(\hgt) &= 75\flux \cdot \big|\{k : 657{,}850 \le k \le \min(\hgt,1{,}313{,}199)\}\big| \\
\mathrm{pow}_{37.5}(\hgt) &= 37.5\flux \cdot \big|\{k : 1{,}313{,}200 \le k \le \min(\hgt,2{,}019{,}999)\}\big| \\
\mathrm{exch}(\hgt) &= 7{,}500{,}000\flux \cdot \indic{\hgt \ge 836{,}274} \\
\mathrm{found}(\hgt) &= 2{,}500{,}000\flux \cdot \indic{\hgt \ge 836{,}994} \\
\mathrm{swap}(\hgt) &= 22{,}000{,}000\flux \cdot \textstyle\sum_{i=0}^{9} \indic{\hgt \ge 837{,}714 + 21{,}600\,i} \\
\mathrm{pon}(\hgt) &= \textstyle\sum_{k=\hpon}^{\hgt} \subsidypon(k) \cdot \indic{\hgt \ge \hpon}
\end{align*}
\end{definition}
\shipped \srccode{flux-roadmap/\allowbreak{}supply/\allowbreak{}emission\_model.py}{321-367}. The premine trigger height
above uses $2$ (the observed mode used for every current- and historical-supply figure in this
section); the literal code branch triggers at height $1$ instead
\shipped \srccode{fluxd/\allowbreak{}src/\allowbreak{}main.cpp}{2820-2822} — the one-block difference re-enters as part of
the discrepancy in \S\ref{sec:mp-discrepancy}. All nine terms are additive; the PoW and PoN
subsidy terms are mutually exclusive by height, and the premine, exchange, foundation and
swap-pool terms are new issuance layered on top of whichever subsidy term is active
\shipped \srccode{flux-roadmap/\allowbreak{}supply/\allowbreak{}emission\_model.py}{469-477}. The tier rewards and dev-fund
remainder of \S\ref{sec:mp-tier-devfund} are a display-only split of $\mathrm{pon}(\hgt)$; they are
never added on top of it \shipped \srccode{flux-roadmap/\allowbreak{}supply/\allowbreak{}emission\_model.py}{358-368}.

\begin{proposition}[No terminal supply]
\label{prop:no-terminal-supply}
$\supply$ does not converge at any height. For every $\hgt \ge \hgt_{20} \defeq \hpon + \redmax\cdot\redint = 23{,}044{,}000$,
$$
\subsidypon(\hgt) = s_{20} = 170{,}207{,}313\sat = 1.70207313\flux
$$
is constant, because the reduction loop in Definition~\ref{def:pon-subsidy} terminates after
$\redmax = 20$ iterations rather than continuing indefinitely
\shipped \srccode{fluxd/\allowbreak{}src/\allowbreak{}main.cpp}{2811-2816}. Consequently $\supply$ grows without bound at the
fixed asymptotic rate
$$
s_{20} \cdot \redint = 1.70207313\flux \times 1{,}051{,}200 = 1{,}789{,}219.274256\ \flux/\mathrm{year}, \quad \text{forever}.
$$
\end{proposition}
\shipped \srccode{flux-roadmap/\allowbreak{}supply/\allowbreak{}emission\_model.py}{514-546}, executed 2026-09-03.
Height $\hgt_{20} = 23{,}044{,}000$ is a direct consequence of the consensus constants above and is
exact; the corresponding calendar date, $\sim$2045-10-20, is a 30-second-spacing extrapolation from
the PoN-activation timestamp anchor, and the model itself validates that extrapolation only against
block $2{,}816{,}820$, about $800{,}000$ blocks after activation, not over horizons this long
\srcdoc{flux-roadmap/supply/emission\_model.py:389-391}.
One PoN year later, at height $24{,}095{,}199$ ($\sim$2046-10-20, same caveat), cumulative supply is
$548{,}043{,}625.860464\flux$ \shipped \srccode{flux-roadmap/\allowbreak{}supply/\allowbreak{}emission\_model.py}{321-367}.
This is not a cap; it is only the point at which the annual emission \emph{rate} has become
constant, per Proposition~\ref{prop:no-terminal-supply}. The team's own public statement is
consistent with this derivation: ``block rewards reduce 10\% a year for twenty years and then
continue at a small fixed amount indefinitely — about 1.79 million FLUX a year, forever''
\srcroadmap{flux-roadmap-public.md:189}.

\textbf{Current supply.} $\supply(2{,}912{,}015) = 429{,}466{,}823.9700\flux$, observed mode,
computed 2026-09-03 \shipped \srccode{flux-roadmap/\allowbreak{}supply/\allowbreak{}emission\_model.py}{321-367}. This is
consistent with the team's own public reconciliation — ``about 428.1 million FLUX, at block
$\sim$2,816,845 \dots reconciles with our emission model to within 0.0005\%'' —
\srcroadmap{flux-roadmap-public.md:185}, an earlier height on the same monotone curve.

\textbf{\maxmoney is not a supply cap.} $\maxmoney = 440{,}000{,}000\flux$ is defined in
\texttt{MoneyRange()} as a per-transaction/output sanity bound, not a chain-wide accumulator
\shipped \srccode{fluxd/\allowbreak{}src/\allowbreak{}amount.h}{22-31}. Cumulative issuance crosses this value at height
$3{,}730{,}295$ ($\sim$2027-06-11) \shipped \srccode{flux-roadmap/\allowbreak{}supply/\allowbreak{}emission\_model.py}{550-561},
and nothing in consensus checks or breaks when it does. The team states this publicly in the same
terms: ``Neither 440 million nor 560 million is a cap enforced by the protocol \dots issuance
passes it in 2027 without anything breaking \dots there is no hard cap today. We would rather say
that than repeat a number'' \srcroadmap{flux-roadmap-public.md:189,191}.

\subsection{Plots}
\label{sec:mp-plots}

\begin{figure}[H]
\centering
\begin{tikzpicture}
\begin{axis}[
    width=\linewidth, height=6cm,
    xlabel={Year},
    ylabel={Cumulative supply $\supply(\hgt)$ [\flux]},
    xmin=2026, xmax=2060,
    grid=major,
]
\addplot[thick, blue, mark=none] table[x index=0, y index=2] {../data/emission.dat};
\draw[densely dashed] ({axis cs:2045.8,0}|-{rel axis cs:0,1}) -- ({axis cs:2045.8,0}|-{rel axis cs:0,0})
    node[pos=0.08, right, font=\scriptsize] {$\hgt=23{,}044{,}000$: 20th reduction};
\end{axis}
\end{tikzpicture}
\caption{Cumulative issued \Flux supply, 2026--2060, mode=observed
\shipped \srccode{flux-roadmap/\allowbreak{}supply/\allowbreak{}emission\_model.py}{321-367}. Data: \texttt{paper/\allowbreak{}data/\allowbreak{}emission.\allowbreak{}dat}
(columns \texttt{year}, \texttt{height}, \texttt{cumulative\_supply}, \texttt{annual\_emission},
\texttt{pow\_emission}, \texttt{pon\_emission}). The dashed line marks the onset of the flat
tail rate, per Proposition~\ref{prop:no-terminal-supply}, read from
\texttt{paper/\allowbreak{}data/\allowbreak{}pon\_reductions.\allowbreak{}dat} row $n=20$.}
\label{fig:cumulative-supply}
\end{figure}

\begin{figure}[H]
\centering
\begin{tikzpicture}
\begin{axis}[
    width=\linewidth, height=6cm,
    xlabel={Year},
    ylabel={Annual emission [\flux/\text{yr}]},
    xmin=2026, xmax=2060,
    grid=major,
]
\addplot[thick, red!70!black, mark=none] table[x index=0, y index=3] {../data/emission.dat};
\end{axis}
\end{tikzpicture}
\caption{Annual \Flux emission by calendar year. \texttt{pow\_emission} is exactly zero for every
plotted year (all postdate \PoN activation) \shipped \srccode{flux-roadmap/\allowbreak{}supply/\allowbreak{}emission\_model.py}{299-318}.
Data: \texttt{paper/\allowbreak{}data/\allowbreak{}emission.\allowbreak{}dat}, column \texttt{annual\_emission}. From calendar year 2046
onward the series is $1{,}789{,}219.274256\flux$ per 365-day year, per
Proposition~\ref{prop:no-terminal-supply}; leap years (2048, 2052, \ldots) carry one more day of
blocks and read $1{,}794{,}121.2449\flux$.}
\label{fig:annual-emission}
\end{figure}
%% verified 2026-09-04: emission.dat header is "# year height cumulative_supply annual_emission
%% pow_emission pon_emission"; the cumulative-supply plots read y index=2 and the annual-emission
%% plot y index=3, both matching their axis labels. pon_reductions.dat is not consumed by any
%% plot. (Build note; not paper text.)

\subsection{A discrepancy in the published model}
\label{sec:mp-discrepancy}

The emission model's own docstring claims that its two computation modes — \texttt{observed}
(what mainnet actually paid) and \texttt{code} (a literal transcription of the current
\texttt{Get\allowbreak{}Block\allowbreak{}Subsidy}) — differ by $0.06\flux$ in total, and calls the choice between them
``irrelevant at publication precision'' \srcdoc{flux-roadmap/supply/emission\_model.py:105-106}.
Executing the model shows an exact, height-independent difference of $150.00000000\flux$ at every
tested height $\ge 5{,}000$, arising from the one-block index shift itself, which lowers the payout of
every one of the $4{,}996$ ramp blocks by $0.03\flux$ and of blocks $2$ and $2{,}500$ by $0.06\flux$
each, where the docstring's figure counts block $2$ alone
\shipped \srccode{flux-roadmap/\allowbreak{}supply/\allowbreak{}emission\_model.py}{207-229}. This is a
finding about the published emission model's own arithmetic, not about \fluxd's consensus code,
which is unaffected. This paper uses \texttt{observed} mode throughout, as the model's own default
for historical and current-supply claims; full derivation and reproduction steps are in
Appendix~C.

\subsection{Baseline for what follows}

The subsidy, tier, dev-fund and fund-payout functions given here are the quantities Section~7
measures Progressive Node Reward parity against, and every issuance, supply or operator-income
figure elsewhere in this paper is denominated against the schedule fixed in this section.

%% Drain this section's float queue before the next begins.
\FloatBarrier
